% KernelForge -- arXiv paper skeleton (target: January 2027).
%
% HOW TO READ THIS FILE
%   * Every \todo{...} is something only the GPU run (or the owner) can fill in.
%     docs/paper/CHECKLIST.md lists them in order and says which claims depend on which result.
%   * Every claim about the repo was checked against the code on 2026-11-15; the file it comes
%     from is given in a "% src:" comment next to it. If the code changes, re-check the claim.
%   * Every claim about other work is limited to what docs/paper/SOURCES.md records.
%   * Build: pdflatex paper && bibtex paper && pdflatex paper && pdflatex paper

\documentclass[11pt]{article}

\usepackage[utf8]{inputenc}
\usepackage[T1]{fontenc}
\usepackage[margin=1in]{geometry}
\usepackage{amsmath}
\usepackage{booktabs}
\usepackage{xcolor}
\usepackage[numbers,sort&compress]{natbib}
\usepackage[hidelinks]{hyperref}
\usepackage{url}

% A visible placeholder for anything that needs a result or an owner decision.
% Search the PDF for "TODO" before posting; there must be none left.
\newcommand{\todo}[1]{\textcolor{red}{\textbf{[TODO:} #1\textbf{]}}}
% Small helper for code identifiers.
\newcommand{\code}[1]{\texttt{#1}}

\title{KernelForge: A Data-Efficient, Explainable Triton Kernel Generator\\for LLM-Serving Operators}
\author{\todo{author name, affiliation}}
\date{\todo{posting date}}

\begin{document}
\maketitle

% =====================================================================================
\begin{abstract}
Large language models can draft GPU kernels, but recent systems that do this well rely on
large training budgets, often reinforcement learning (RL) on 8B--32B models, and report a
kernel and a benchmark number without saying why the kernel is faster.
We study a deliberately small setting: supervised QLoRA fine-tuning of a
1.5B-parameter code model on a few hundred harness-verified examples, on a single 16\,GB
GPU, with no RL. The scope is narrow on purpose. It covers five operator families that matter
for LLM inference serving: quantized matmul, prefill attention, KV-cache decode, normalization
and rotary embeddings. The model outputs a Triton kernel together with a one-to-three-sentence
explanation of the optimization.
Every training example, every evaluation score, and every kernel shown in our demo is gated by
the same subprocess verification harness. It follows KernelBench's \code{Model}/\code{ModelNew}
contract and checks correctness against the PyTorch reference and measures speedup.
We evaluate on held-out templates, so no test operator appears in training at any shape.
On this split, fine-tuning changes kernel correctness from \todo{X\%} to \todo{Y\%}
\todo{and mean speedup of correct kernels from A$\times$ to B$\times$}. Training-free best-of-4
sampling with one round of harness-feedback repair reaches \todo{Z\%}. On a manual 0/1/2
rubric, \todo{N} of \todo{M} sampled explanations name the specific technique.
\todo{One sentence on the weakest category, stated plainly.}
We do not claim to match larger RL-trained systems. We report what this budget achieves and
where it fails.
\end{abstract}

% =====================================================================================
\section{Introduction}

A few operators account for much of the GPU time in LLM serving: weight-quantized
matrix multiplication, attention during prefill, attention against the KV cache during decode,
normalization layers, and rotary position embeddings. Hand-written Triton~\cite{tillet2019triton}
kernels for operators like these give real gains. Liger Kernel~\cite{hsu2024liger} reports
about 20\% higher training throughput and about 60\% lower memory from fusion and chunking,
and FlashAttention~\cite{dao2022flashattention,dao2024flashattention2} restructures attention around
on-chip memory. Writing such kernels takes expert time, which makes them an appealing
target for LLM code generation.

Recent work has made LLM kernel generation measurable and competitive.
KernelBench~\cite{ouyang2025kernelbench} and TritonBench~\cite{li2025tritonbench} provide benchmarks.
KernelLLM~\cite{fisches2025kernelllm} fine-tunes an 8B model on roughly 25k PyTorch--Triton pairs.
Kevin~\cite{baronio2025kevin}, AutoTriton~\cite{li2025autotriton}, TritonRL~\cite{woo2025tritonrl} and
Dr.~Kernel~\cite{liu2026drkernel} train 8B--32B models with reinforcement learning, usually after a
supervised stage. These systems are better resourced than this project, and we do not try to
beat them on broad coverage or raw speedup.

We ask a narrower question instead: \emph{with a small curated supervised dataset, a 1.5B model,
no RL, and one free-tier 16\,GB GPU, how good can a kernel generator get on a narrow,
serving-relevant set of operators, and can it say why its kernel is faster?}

\paragraph{Contributions.}
\begin{enumerate}
  \item \textbf{A harness-gated dataset for serving operators.} We built 15 hand-written Triton
        templates across five categories, instantiated at 20 shapes each (300 candidates). An
        example enters the dataset only if a subprocess harness confirms that it matches the
        PyTorch reference (\S\ref{sec:dataset}).
  \item \textbf{An explanation output and its evaluation.} The training target contains a short
        rationale next to the kernel. We score generated explanations with a manual rubric
        (\S\ref{sec:explanations}). None of the works we reviewed evaluates a
        natural-language explanation as an output \todo{re-check against full papers before posting; see UNVERIFIED.md}.
  \item \textbf{One verification function everywhere.} The same function filters training data,
        scores the baseline and the fine-tuned model, and drives an inference-time best-of-$N$ +
        repair search. So ``correct'' means the same thing in every number we report
        (\S\ref{sec:harness}, \S\ref{sec:search}).
  \item \textbf{An honest small-budget reference point.} We report category-level results on a
        held-out-template split, including failures and weak categories (\S\ref{sec:results}).
\end{enumerate}

% =====================================================================================
\section{Related Work}
\label{sec:related}

\paragraph{Benchmarks.}
KernelBench~\cite{ouyang2025kernelbench} collects 250 PyTorch workloads. It asks a model for a
faster \code{ModelNew} that matches a reference \code{Model}, and it introduces $\mathrm{fast}_p$,
the fraction of generated kernels that are correct \emph{and} faster than the reference by more
than a factor $p$. Its authors find that frontier reasoning models match the PyTorch baseline in
under 20\% of cases out of the box, and that execution and profiling feedback helps. We adopt
its \code{Model}/\code{ModelNew} contract unchanged. TritonBench~\cite{li2025tritonbench} targets
Triton specifically, with 184 real-world GitHub operators (TritonBench-G) and a PyTorch-aligned
set (TritonBench-T). The best execution accuracy it reports is 23.91\% on G and 53.01\% on T, and
the best speedups among correct operators are 1.56$\times$ and 1.91$\times$. Both benchmarks cover
general operators. We cover five serving-operator families only, so our numbers are \emph{not}
comparable to either benchmark's leaderboard.

\paragraph{Supervised kernel generators.}
KernelLLM~\cite{fisches2025kernelllm} is the closest prior work to our setup. It is an 8B model
based on Llama~3.1 Instruct, fine-tuned with supervised learning on about 25k PyTorch--Triton
pairs plus synthetic samples, using 16 GPUs for about 12 hours. It reports 20.2 pass@1 on
KernelBench-Triton Level~1. Our setting differs in four ways: about two orders of magnitude
less data (\todo{final verified count} examples versus about 25k), a model about five times
smaller, a serving-operator scope instead of general PyTorch modules, and an explanation output.

\paragraph{Reinforcement-learning kernel generators.}
Kevin~\cite{baronio2025kevin} trains QwQ-32B with multi-turn RL to write and refine
\emph{CUDA} kernels. On its evaluation it raises correctness from 56\% to 82\% and mean speedup
from 0.53$\times$ to 1.10$\times$, and it finds that serial refinement scales better than
parallel sampling. AutoTriton~\cite{li2025autotriton} (8B) and TritonRL~\cite{woo2025tritonrl}
(8B) both run supervised fine-tuning followed by RL with execution-based rewards. TritonRL
adds layered verification and a hierarchical reward to counter reward hacking.
Dr.~Kernel~\cite{liu2026drkernel} (14B) trains with multi-turn RL in a distributed GPU environment
with reward-hacking checks. We use no RL. Our repair loop (\S\ref{sec:search}) borrows Kevin's
multi-turn feedback idea but runs it only at inference time, with an untrained base model or
an SFT adapter. We do not claim it reaches what trained refinement reaches.

\paragraph{Training-free agents.}
GEAK~\cite{wang2025geak} generates Triton kernels for AMD GPUs with an agent loop of
generation, evaluation, reflection and optimization. Sakana AI's agentic pipeline~\cite{lange2025robustkbench}
translates PyTorch to CUDA with evolutionary search and LLM-based verifiers. It also argues
that existing kernel benchmarks have exploitable loopholes, which is one reason we keep a
single conservative harness (\S\ref{sec:harness}). Both rely on frontier models at test time.
Our search uses a 1.5B model and a budget of at most four samples and two repair rounds. More
generally, feeding execution results back to a code model goes back to Self-Debugging~\cite{chen2023selfdebug}
and Reflexion~\cite{shinn2023reflexion}.

\paragraph{Kernels we draw on.}
Our attention templates follow FlashAttention~\cite{dao2022flashattention,dao2024flashattention2}:
query-block programs, KV tiles, and online softmax~\cite{milakov2018online}. Our decode templates
use a tiled online-softmax loop over the KV cache. Unlike Flash-Decoding~\cite{dao2023flashdecoding},
they do \emph{not} split the KV sequence across programs. GQA~\cite{ainslie2023gqa},
RMSNorm~\cite{zhang2019rmsnorm} and rotary embeddings~\cite{su2021roformer} define the reference
operators.

\paragraph{Positioning.}
Compared with the systems above, this work (i)~has a narrower scope (five serving-operator families);
(ii)~uses no RL; (iii)~uses supervised QLoRA~\cite{dettmers2023qlora} on a 1.5B model;
(iv)~uses a few hundred examples, every one of which passed the harness; and (v)~generates and evaluates an
explanation. These are differences in setting, not claims of superiority.

% =====================================================================================
\section{Method}
\label{sec:method}

\subsection{Task}
The input is a PyTorch reference implementation of a serving operator plus a one-paragraph
description. The output is (a)~a self-contained Python module that defines Triton kernels and a
\code{ModelNew} class, and (b)~a 1--3 sentence explanation of the optimization technique and why
it helps at serving scale. % src: kernelforge/prompts.py OUTPUT_FORMAT_INSTRUCTION
The model must answer in two labelled sections, \code{kernel:} and \code{explanation:}, which a
parser extracts. % src: kernelforge/parsing.py
The prompt also states the contract: import only \code{torch}, \code{triton},
\code{triton.language} and \code{math}; give \code{ModelNew} the same constructor and \code{forward}
signature as \code{Model}; allocate outputs on the input's device; do not redefine
\code{get\_inputs}/\code{get\_init\_inputs}. % src: kernelforge/prompts.py KERNEL_CONTRACT

\subsection{Operator categories}
We fixed five categories with three templates each (Table~\ref{tab:categories}). All references
are fp16. % src: data/templates/*.py (all entries dtype float16)
The choice is driven by serving, not training. Decode attention and quantized matmul dominate
token-generation cost, and norms and RoPE run in every layer of every forward pass.
\todo{optional: cite a serving-cost profile if you want to make this quantitative; none is cited now.}

\begin{table}[t]
\centering
\small
\begin{tabular}{@{}ll@{}}
\toprule
Category & Templates \\
\midrule
\code{quantized\_matmul} & W8A16 per-tensor scale; W8A16 per-channel scale; W8A8 \\
\code{attention} (prefill) & causal self-attention; non-causal cross-attention; GQA causal \\
\code{kv\_cache} (decode) & decode attention (MHA); decode attention (GQA); KV-cache write \\
\code{norm} & RMSNorm; LayerNorm; fused residual-add + RMSNorm \\
\code{rope} & rotate-half; interleaved pairs; fused Q+K \\
\bottomrule
\end{tabular}
\caption{The 15 templates. Each is instantiated at 20 shapes, giving 300 candidates.}
\label{tab:categories}
\end{table}

\subsection{Template library and harness-gated dataset}
\label{sec:dataset}
Writing 300 genuinely different correct kernels by hand is not feasible for one person, so each
category has three hand-written templates crossed with a grid of 20 shape configurations.
% src: data/templates/__init__.py GENERATORS; data/build_dataset.py
Every entry stores the reference, the kernel, the explanation, the dtype, a tolerance
($\mathrm{atol}=\mathrm{rtol}=10^{-2}$ for fp16), the test shapes, and a verification record.
% src: data/SCHEMA.md; data/templates/common.py FP16_TOLERANCE
Each candidate is run through the harness on the GPU. Only passing entries enter the dataset, and
the rest are kept with their failure reason. Out-of-memory failures are counted separately from
incorrect kernels, so a memory limit on the 16\,GB GPU is not mistaken for a broken template.
% src: data/build_dataset.py; MASTER_PLAN.md Part III

Before any GPU time, every template was run on a small shape through Triton's CPU interpreter
in a Linux container. All 15 pass. % src: data/raw/interpret_report.json (15/15)
This check found real bugs, including attention references that failed on a mixed-dtype matmul
and a W8A8 reference that used an int32 matmul. It checks correctness only and gives no timing.
None of the results below rely on it.

Two properties of this dataset matter for interpreting the results:
\begin{itemize}
  \item \textbf{Within a template, entries differ only in shape constants.} The kernel logic is the
        same, so the effective number of distinct kernels is 15, not 300.
  \item \textbf{The explanation is fixed per template.} Every template except the KV-cache write
        has a single explanation string shared by all 20 shapes, and the KV-cache write has four.
        % src: computed from data/templates generators on 2026-11-15
        The training set therefore contains about a dozen distinct explanations. Our explanation
        evaluation (\S\ref{sec:explanations}) is designed around this. It scores explanations on
        held-out templates, where copying a training explanation would be visibly wrong.
\end{itemize}

\subsection{Verification harness}
\label{sec:harness}
The harness follows KernelBench's contract~\cite{ouyang2025kernelbench}. The reference module defines
\code{Model}, \code{get\_inputs()} and \code{get\_init\_inputs()}, and the candidate defines
\code{ModelNew} with the same signatures. % src: verification/verify_kernel.py verify_model_kernel
Both are written to temporary files and executed in a separate Python subprocess with a
wall-clock timeout. A hang becomes \code{timeout} and a hard crash becomes \code{crash}, so neither
takes down the calling loop. % src: verification/verify_kernel.py, sandbox_runner.py
Inside the subprocess:
\begin{itemize}
  \item The RNG is seeded immediately before \emph{each} of \code{Model(*init)} and
        \code{ModelNew(*init)}, so operators with random internal state (e.g.\ quantized weights)
        compare against identical state.
  \item Reference and candidate get separately constructed input tensors from the same seed. This
        catches a candidate that mutates its inputs in place.
  \item Outputs (including tuples) are cast to fp32 and compared with
        \code{torch.allclose} at the entry's tolerance, with \code{equal\_nan=False}.
  \item On a GPU, each side is timed with CUDA events after untimed warm-up iterations, and the
        median is reported. Speedup is $t_{\mathrm{ref}}/t_{\mathrm{cand}}$, where the
        reference is PyTorch eager. Dataset building uses 10 warm-up and 30 timed iterations, and
        evaluation uses 5 and 20. % src: data/build_dataset.py, evaluation/baseline_eval.py defaults
  \item A result \emph{passes} only if it ran (\code{status} \code{ok}) and is correct.
\end{itemize}
When outputs differ, the harness describes the difference (shape, output count, NaN count, or the
number of elements outside tolerance and the location of the largest error). The repair loop
feeds this description back to the model.

\paragraph{Interpreter mode versus GPU mode.}
With \code{interpret=True}, the subprocess runs with \code{TRITON\_INTERPRET=1} on CPU tensors. It
returns correctness only, marked as interpreted, with no timings. We use it for local
development and the hosted demo. \emph{All numbers in this paper come from GPU mode} \todo{confirm
when filling results}.

\paragraph{Scope of the sandbox.}
The subprocess has rlimits and a scrubbed environment, and the public demo adds an AST import
policy. It is not a security boundary in the sense of a VM or seccomp. % src: DOCUMENTATION.md "Safety scope"
This matters for a public demo, not for the results.

\subsection{Train/validation/test split by template}
A random shape-level split would put the same kernel, with different constants, in both train and
test, and test accuracy would then measure memorization. We instead hold out one \emph{whole
template} per category for test. From the remaining templates, 10\% of entries go to validation.
% src: kernelforge/dataset.py split_by_template (seed 42)
If all 300 candidates pass on the GPU, seed 42 holds out GQA causal prefill attention, GQA decode
attention, fused residual-add + RMSNorm, W8A16 per-channel matmul and interleaved RoPE, and the
split is 180/20/100. \todo{replace with the actual held-out templates and counts from
\code{split\_info.json}; the split is computed on the verified set, so failures can change it.}

Each held-out template has a close sibling in training: GQA versus MHA/causal attention,
per-channel versus per-tensor scaling, RMSNorm with versus without a fused residual. The test set
therefore measures generalization to \emph{variants} of seen operator families, not to new
families. We state this as a property of the benchmark, not a weakness we hide.

\subsection{Training format}
Each example is a chat transcript: a fixed system prompt, a user turn with the description and
reference followed by the output contract, and an assistant turn with
\code{kernel:}\,\dots\,\code{explanation:}\,\dots. It is rendered with the base model's chat template.
% src: training/finetune.py _format_chat_example; kernelforge/prompts.py
The loss is the standard causal-LM loss over the whole rendered sequence, with no masking of the
prompt tokens. % src: training/finetune.py (DataCollatorForLanguageModeling / SFTTrainer text field)
\todo{owner decision: keep, or switch to completion-only loss and re-run; state which.}

\subsection{Inference: verifier-guided search and repair}
\label{sec:search}
The search is pure orchestration around the harness. % src: kernelforge/search.py
\begin{enumerate}
  \item \textbf{Sample.} Draw $N\in\{1,\dots,4\}$ completions in one batched generation. With
        $N>1$ the temperature is raised to at least 0.7, so the samples differ.
  \item \textbf{Verify.} Parse each completion and run each parseable kernel through the harness.
        Identical kernels are verified once.
  \item \textbf{Select.} Among passing candidates, pick the one with the highest measured speedup.
  \item \textbf{Repair.} If none pass, take the most useful failure (wrong numerics, then a harness
        error, then unparseable output) and send it back to the model as the next turn: the original
        prompt, that attempt, and the harness feedback (the mismatch description or the last 25
        traceback lines). Do this for up to two rounds. Only the latest attempt is kept in context.
\end{enumerate}
For general-purpose (non-fine-tuned) models, the demo can add a solved example from a
\emph{different} template in the same category, a short list of Triton~3.x pitfalls, and
pattern-matched repair hints. % src: kernelforge/search.py, kernelforge/prompts.py TRITON_NOTES, repair_hints
\todo{state exactly which of these were on for each row of Table~\ref{tab:search}. The current
evaluation scripts use a zero-shot prompt with no example and no notes.}

% =====================================================================================
\section{Experimental Setup}
\label{sec:setup}

\paragraph{Model.} Qwen2.5-Coder-1.5B-Instruct~\cite{hui2024qwen25coder}.
% src: training/configs/qwen2.5-coder-1.5b.json

\paragraph{Fine-tuning.} QLoRA~\cite{dettmers2023qlora}: the base model is loaded in 4-bit NF4 with
double quantization and fp16 compute. LoRA~\cite{hu2022lora} adapters have rank $r=16$,
$\alpha=32$, dropout 0.05, on the \code{q,k,v,o,gate,up,down} projections. Training uses
learning rate $2\times10^{-4}$, 3 epochs, per-device batch size 1 with 8 gradient-accumulation
steps (effective batch 8), warm-up ratio 0.05, weight decay 0.01, maximum sequence length 8192,
seed 42 and fp16 mixed precision. % src: training/configs/finetune_default.json; training/finetune.py
The preferred backend is Unsloth~\cite{unsloth2023}. Where it conflicts with the preinstalled
PyTorch, we fall back to transformers + PEFT + bitsandbytes with the same hyperparameters and the
\code{paged\_adamw\_8bit} optimizer. \todo{which backend actually ran; optimizer if Unsloth.}

\paragraph{Hardware.} One NVIDIA T4 (16\,GB) in a free Kaggle session, for dataset verification,
training and evaluation. \todo{driver/CUDA/PyTorch/Triton versions from the session's env
report; wall-clock time for each stage; peak memory during training.}

\paragraph{Decoding.} Temperature 0.2, top-$p$ 0.95, one sample per prompt for pass@1, and a cap of
\todo{1536} new tokens. % src: training/configs/qwen2.5-coder-1.5b.json; scripts/gpu_session.py --max-new-tokens
\todo{check the cap does not truncate the longest reference kernels (about 4.3k characters). If
it does, raise it and re-run.}

\paragraph{Metrics.} For each category we report:
\begin{itemize}
  \item \textbf{\% parsed:} the fraction of outputs from which a kernel could be extracted.
  \item \textbf{\% correct:} the fraction of \emph{all} test prompts whose kernel passed the harness.
  \item \textbf{Mean speedup:} the arithmetic mean of $t_{\mathrm{ref}}/t_{\mathrm{cand}}$ over
        \emph{passing} kernels only. % src: evaluation/common.py CategoryMetrics.avg_speedup
        Because it is conditional on correctness, a model that solves only easy cases can show a higher
        mean speedup. We therefore also report $\mathrm{fast}_1$ (correct \emph{and} faster than eager),
        following KernelBench~\cite{ouyang2025kernelbench}. \todo{fast\_1 and geometric mean are not computed
        by the current scripts; add them or drop this sentence.}
\end{itemize}
Test prompts in a category share one held-out template, so the 20 shapes are not independent
samples. We report counts next to percentages and give Wilson 95\% intervals.
\todo{n per category: the default session evaluates 25 prompts in total (5 per category). Decide
whether to evaluate all held-out prompts.}

\paragraph{Baselines.} The same base model without an adapter, with the same prompt, decoding
and harness. \todo{optional: a second base model (e.g.\ DeepSeek-Coder or StarCoder2) if time allows.}

% =====================================================================================
\section{Results}
\label{sec:results}

\todo{Write this section only after the tables are filled. Lead with the weakest category.}

\begin{table}[t]
\centering
\small
\begin{tabular}{@{}lrrrrr@{}}
\toprule
Category & Generated & GPU-verified & Incorrect & OOM & Other error \\
\midrule
Quantized matmul & 60 & \todo{} & \todo{} & \todo{} & \todo{} \\
Attention        & 60 & \todo{} & \todo{} & \todo{} & \todo{} \\
KV-cache         & 60 & \todo{} & \todo{} & \todo{} & \todo{} \\
Norm             & 60 & \todo{} & \todo{} & \todo{} & \todo{} \\
RoPE             & 60 & \todo{} & \todo{} & \todo{} & \todo{} \\
\midrule
Total            & 300 & \todo{} & \todo{} & \todo{} & \todo{} \\
\bottomrule
\end{tabular}
\caption{Dataset construction on one T4. Source: \code{data/raw/generation\_report.json}.
\todo{add median speedup of the hand-written templates per category; this is the ceiling the
model is trained toward.}}
\label{tab:dataset}
\end{table}

\begin{table}[t]
\centering
\small
\setlength{\tabcolsep}{4pt}
\begin{tabular}{@{}lrrrrrr@{}}
\toprule
 & \multicolumn{2}{c}{Base} & \multicolumn{2}{c}{Fine-tuned} & FT best-of-4 + repair & \\
\cmidrule(lr){2-3}\cmidrule(lr){4-5}\cmidrule(lr){6-6}
Category & \% correct & mean speedup & \% correct & mean speedup & \% correct & $n$ \\
\midrule
Quantized matmul & \todo{} & \todo{} & \todo{} & \todo{} & \todo{} & \todo{} \\
Attention        & \todo{} & \todo{} & \todo{} & \todo{} & \todo{} & \todo{} \\
KV-cache         & \todo{} & \todo{} & \todo{} & \todo{} & \todo{} & \todo{} \\
Norm             & \todo{} & \todo{} & \todo{} & \todo{} & \todo{} & \todo{} \\
RoPE             & \todo{} & \todo{} & \todo{} & \todo{} & \todo{} & \todo{} \\
\bottomrule
\end{tabular}
\caption{Held-out-template test set, GPU harness. Mean speedup is over passing kernels only
(``--'' if none passed). \todo{add \% parsed and Wilson intervals, or put them in an appendix table.}}
\label{tab:main}
\end{table}

\begin{table}[t]
\centering
\small
\begin{tabular}{@{}lrrr@{}}
\toprule
Model & pass@1 & best-of-4 & best-of-4 + 1 repair \\
\midrule
Base       & \todo{} & \todo{} & \todo{} \\
Fine-tuned & \todo{} & \todo{} & \todo{} \\
\bottomrule
\end{tabular}
\caption{Training-free search, \% correct over all test prompts. Extra cost per prompt:
\todo{mean generated tokens and harness runs per column}. \todo{the evaluation scripts do not run
the search yet; see CHECKLIST.md.}}
\label{tab:search}
\end{table}

\begin{table}[t]
\centering
\small
\begin{tabular}{@{}lrrrr@{}}
\toprule
 & Train loss (final) & Val loss (best) & Wall clock & Peak GPU memory \\
\midrule
QLoRA run & \todo{} & \todo{} & \todo{} & \todo{} \\
\bottomrule
\end{tabular}
\caption{Training run. Source: training logs in the session output archive.}
\label{tab:training}
\end{table}

\paragraph{Failure modes.} \todo{Categorize failing outputs: unparseable, harness error (e.g.\
compile or shape error), wrong numerics, timeout. Give counts per model. Show 3--5 concrete examples,
good and bad, in an appendix.}

% =====================================================================================
\section{Explanation Evaluation}
\label{sec:explanations}

Explanation quality has no reliable automatic metric, so we score it by hand.
\todo{20--30} outputs are sampled at random (seed \todo{}) from each evaluation report. Each
is shown next to the reference explanation for its template and scored with a three-level rubric.
% src: evaluation/explanation_review.py (default --sample 25)
\begin{itemize}
  \item[\textbf{2}] correctly identifies the specific optimization technique;
  \item[\textbf{1}] partially correct, or generic but pointing the right way;
  \item[\textbf{0}] wrong, missing, or pure boilerplate.
\end{itemize}
We also check whether an explanation describes what the generated kernel \emph{actually does}, not
only what the reference does. \todo{decide whether to score this as a separate yes/no column.}
Because training explanations are fixed per template (\S\ref{sec:dataset}), we also flag
explanations copied verbatim, or nearly so, from a training template. On a held-out template such a
copy usually describes the wrong operator.

\begin{table}[t]
\centering
\small
\begin{tabular}{@{}lrrrr@{}}
\toprule
Model & Score 2 & Score 1 & Score 0 & Copied from a training template \\
\midrule
Base       & \todo{} & \todo{} & \todo{} & -- \\
Fine-tuned & \todo{} & \todo{} & \todo{} & \todo{} \\
\bottomrule
\end{tabular}
\caption{Manual explanation rubric on \todo{n} samples per model.}
\label{tab:rubric}
\end{table}

\paragraph{Rater.} \todo{One rater (the author), not blinded to model identity, or blinded by
shuffling base and fine-tuned outputs? State which. A second rater on a subset would allow an
agreement figure.}

% =====================================================================================
\section{Limitations}
\label{sec:limitations}
\begin{itemize}
  \item \textbf{Scale.} A 1.5B model and a dataset of a few hundred entries built from 15
        templates, compared with 8B--32B systems trained on tens of thousands of examples or with RL.
  \item \textbf{Coverage.} Five narrow operator families, fp16 only, with no backward passes. The test
        set holds out variants of seen families, not new families.
  \item \textbf{Low diversity.} Shape variants of one template share their kernel logic and their
        explanation, so per-category results rest on one held-out operator each.
  \item \textbf{No RL.} The search improves results at inference time only, at the cost of extra
        generations and harness runs per request.
  \item \textbf{Explanations are scored by hand.} The rubric is coarse, the rater is the author
        \todo{adjust if blinded or double-rated}, and there is no automated metric.
  \item \textbf{Hardware.} All timings come from one T4. Speedups over PyTorch eager on a Turing GPU
        may not transfer to newer GPUs. Timings are medians of 20--30 runs on a shared cloud VM.
  \item \textbf{Interpreter.} The CPU interpreter checks correctness on small shapes only, and we never
        report it as GPU verification.
  \item \textbf{Sandbox.} Subprocess isolation with rlimits is suitable for research and a short-lived
        demo, not for running untrusted code in production.
  \item \textbf{Baseline strength.} \todo{if the base model is evaluated zero-shot, say that a few-shot
        prompt might narrow the gap, and report it if run.}
\end{itemize}

% =====================================================================================
\section{Future Work}
\begin{itemize}
  \item \textbf{Rejection-sampling fine-tuning.} Run the search over the training prompts, keep
        kernels that pass and beat a speedup threshold, and fine-tune on them. This is a lightweight
        step toward learning from the verifier without full RL.
  \item \textbf{Autotuning.} Teach \code{@triton.autotune} configurations instead of fixed block sizes.
  \item \textbf{More fused operators.} Fusion is where most speedups come from, and more templates
        per category would raise diversity.
  \item \textbf{Split-KV decode.} Add a Flash-Decoding-style~\cite{dao2023flashdecoding} decode template
        for long contexts.
  \item \textbf{Serving the generator.} Merge the adapter, quantize it, and serve it with a production
        inference engine.
\end{itemize}

% =====================================================================================
\section{Reproducibility}
Code, templates, the harness and all scripts are at \todo{repository URL and commit hash}.
The whole GPU pipeline (environment report, harness tests on the GPU, dataset build, baseline,
fine-tuning and fine-tuned evaluation) runs from one notebook
(\code{notebooks/kaggle\_gpu\_session.ipynb}, which calls \code{scripts/gpu\_session.py}) on a free
Kaggle T4 and writes every artifact to one archive. % src: scripts/gpu_session.py
We release \todo{the verified dataset, the adapter, the evaluation reports, the rubric scores and
the session archive}. The CPU interpreter check can be reproduced without a GPU:
\begin{verbatim}
docker build -f docker/verify.Dockerfile -t kernelforge-verify .
docker run --rm -v "$PWD":/work kernelforge-verify \
    python data/build_dataset.py --interpret --smoke --jobs 8
\end{verbatim}
Seeds: 42 for the split and for training. \todo{report the run-to-run variance of speedups if the
timing step is repeated; otherwise state that it was not measured.}

\bibliographystyle{plainnat}
\bibliography{references}

\end{document}
